% \iffalse meta-comment
%
% Copyright (C) 2019-2023 by Jan Hajer
% -----------------------------------
%
% This file may be distributed and/or modified under the
% conditions of the LaTeX Project Public License, either version 1.3c
% of this license or (at your option) any later version.
% The latest version of this license is in:
%
% http://www.latex-project.org/lppl.txt
%
% and version 1.3c or later is part of all distributions of LaTeX
% version 2005/12/01 or later.
%
% \fi
%
% \iffalse

%<package>\NeedsTeXFormat{LaTeX2e}[2021-11-15]
%<package>\ProvidesPackage{hep-title}[2026/11/13 v1.6 Extended title data: subtitles, series, preprints, affiliations, ORCID, keywords and abstracts]
%<documentation>\ProvidesFile{hep-title-documentation.tex}[2026/11/13 v1.6 hep-title documentation]
%
%<*documentation>

\RequirePackage[l2tabu, orthodox]{nag}

\documentclass{ltxdoc}

\renewcommand*\theCodelineNo{\rmfamily\tstyle\footnotesize\arabic{CodelineNo}}
\AtBeginEnvironment{macrocode}{\renewcommand*{\ttdefault}{clmt}}
\renewcommand*{\MacroFont}{\codestyle}
\AtBeginDocument{\DeleteShortVerb{\|}}
\AtBeginDocument{\MakeShortVerb{\"}}
\EnableCrossrefs
\CodelineIndex
\RecordChanges

\usepackage{hologo}

\usepackage[parskip,oldstyle,font=10pt]{hep-paper}
\bibliography{bibliography}

\acronym{PDF}{portable document format}

%</documentation>

%<*driver>
\expandafter\newif\csname ifshort\endcsname
\shortfalse
\begin{document}
\DocInput{hep-title-implementation.dtx}
\end{document}
%</driver>
%
% \fi
%
% \CheckSum{1286}
%
% \CharacterTable
%  {Upper-case    \A\B\C\D\E\F\G\H\I\J\K\L\M\N\O\P\Q\R\S\T\U\V\W\X\Y\Z
%   Lower-case    \a\b\c\d\e\f\g\h\i\j\k\l\m\n\o\p\q\r\s\t\u\v\w\x\y\z
%   Digits        \0\1\2\3\4\5\6\7\8\9
%   Exclamation   \!     Double quote  \"     Hash (number) \#
%   Dollar        \$     Percent       \%     Ampersand     \&
%   Acute accent  \'     Left paren    \(     Right paren   \)
%   Asterisk      \*     Plus          \+     Comma         \,
%   Minus         \-     Point         \.     Solidus       \/
%   Colon         \:     Semicolon     \;     Less than     \<
%   Equals        \=     Greater than  \>     Question mark \?
%   Commercial at \@     Left bracket  \[     Backslash     \\
%   Right bracket \]     Circumflex    \^     Underscore    \_
%   Grave accent  \`     Left brace    \{     Vertical bar  \|
%   Right brace   \}     Tilde         \~}
%
% \changes{v1.0}{2021/09/01}{First release on CTAN: series, subtitle, preprint placement, authors, affiliations, emails, editors, endorsers, keywords, and abstracts.}
% \changes{v1.1}{2022/11/01}{Added keyword metadata, printed affiliations only when present, and fixed affiliation references.}
% \changes{v1.2}{2023/07/01}{Refreshed the documentation.}
% \changes{v1.3}{2024/11/01}{Released without changes.}
% \changes{v1.4}{2025/09/01}{Added reprint blocks, \cs{doi}, \cs{arxiv}, multiple preprint and reprint lines, and \cs{placereprint}, set preprints ragged, and applied starred command definitions.}
% \changes{v1.5}{2026/08/01}{Raised the LaTeX format requirement and expanded the documentation and README.}
% \changes{v1.6}{2026/11/13}{Measured the preprint and reprint boxes into lengths before placing them, because \cs{put} no longer accepts \software{calc} expressions, and dropped the \software{picture} package, stored the authors unexpanded and set the \PDF author, title, and keywords at \cs{maketitle} without overwriting user settings, dropping labelled footnotes and turning line breaks into spaces, made \cs{author} compatible with the \software{hyperref} option \texttt{pdfusetitle}, fixed the separator between two endorsers, removed spurious spaces before the author line, the \cs{orcmail} mark, and in \cs{reprint}, restricted the ragged affiliation margins to the affiliation block, made \cs{email} usable in \PDF strings, generated the editor and endorser macros from a common definition, added \cs{orcid}, which shares the footnote with \cs{email}, allowed \cs{email}, \cs{orcid}, \cs{orcmail}, and \cs{thanks} directly after \cs{author}, \cs{editor}, and \cs{endorser}, printed all authors if there are no affiliations, kept editors out of the author blocks, and removed a spurious space after the \cs{thanks} mark. Accepted the abstract also before \cs{maketitle}, in the preamble and in the command form \cs{abstract}\marg{text}, and placed it after the title, and broke affiliations preferably after commas, but only as long as this needs no additional line.}
%
% \ifshort
%<*documentation>
% \fi
%
\GetFileInfo{hep-title.sty}

\title{The \software{hep-title} package\thanks{This document corresponds to \software{hep-title}~\fileversion.}}
\subtitle{Extended title data: subtitles, series, preprints, affiliations, ORCID, keywords and abstracts}
\author{Jan Hajer \email{jan.hajer@tecnico.ulisboa.pt}}
\date{\filedate}

% \ifshort
\begin{document}
% \fi

\newgeometry{vscale=.8, vmarginratio=3:4, includeheadfoot, left=11em, marginparwidth=4.6cm, marginparsep=3mm, right=7em}

\maketitle

\begin{abstract}
The \software{hep-title} package extends the title macros of the standard classes.
It adds series and subtitle metadata, preprint and reprint blocks, DOI and arXiv helpers, author, editor, endorser, affiliation, email and ORCID helpers, date and abstract formatting, keywords, and two-column title-aware abstracts.
\end{abstract}

To use the \software{hep-title} package include it with "\usepackage{hep-title}".

Most macros are intended for the use in the preamble.
They are implemented using the \software{titling} \cite{titling} and \software{authblk} \cite{authblk} packages.

\section{Macros}

The \software{orcidlink} package \cite{orcidlink}, which is loaded by \software{hep-title}, always loads the \software{hyperref} package \cite{hyperref}.
Hence, options for \software{hyperref} have to be given with "\hypersetup"\marg{options} or with "\PassOptionsToPackage"\marg{options}"{hyperref}" before \software{hep-title} is loaded; loading "\usepackage"\oarg{options}"{hyperref}" afterwards results in an option clash.
"\maketitle" sets the \PDF title, author, and keywords according to the "\title"\marg{text}, "\subtitle"\marg{text}, "\author"\marg{text}, and "\keywords"\marg{keywords} information.
Footnotes, "\thanks", "\email", "\orcid", "\orcmail", and "\orcidlink" are dropped, line breaks become spaces, and "\and" becomes a comma.
Entries that have been set with "\hypersetup" are kept.
In other \PDF strings, such as bookmarks, "\email" is replaced by the email address and "\orcid" is dropped.

\subsection{Title}

\DescribeMacro{\preseries}
\DescribeMacro{\series}
\DescribeMacro{\postseries}
\DescribeMacro{\title}
\DescribeMacro{\presubtitle}
\DescribeMacro{\subtitle}
\DescribeMacro{\postsubtitle}
\DescribeMacro{\seriesfont}
\DescribeMacro{\titlefont}
\DescribeMacro{\subtitlefont}
The "\series"\marg{series} places a series title above the usual title generated via "\title"\marg{title}.
The "\subtitle"\marg{subtitle} macro places a subtitle below the usual title.
The "\preseries", "\postseries", "\presubtitle", and "\postsubtitle" hooks adjust the surrounding material.
The fonts and their size can be adjusted using the "\seriesfont"\marg{font}, "\titlefont"\marg{font}, and "\subtitlefont"\marg{font} macros.

\subsection{Authors}

\DescribeMacro{\author}
\DescribeMacro{\authorfont}
\DescribeMacro{\authortitle}
\DescribeMacro{\authortitlefont}
\DescribeMacro{\preauthortitle}
\DescribeMacro{\postauthortitle}
\DescribeMacro{\email}
\DescribeMacro{\orcid}
\DescribeMacro{\orcmail}
\DescribeMacro{\affiliation}
\DescribeMacro{\affiliationfont}
In order to facilitate multiple authors with different affiliations the \software{authblk} package \cite{authblk} is extended.
The author macro is extended to take an affiliation label "\author"\oarg{label}\marg{name}.
The affiliation macro takes the corresponding label "\affiliation"\oarg{label}\marg{institution}.
Additionally, the "\email"\marg{email} macro places a link containing the email address in typewriter font into a footnote, and the "\orcid"\marg{orcid} macro places the ORCID logo and identifier, linked to the ORCID record.
Consecutive "\email" and "\orcid" macros share one footnote, so that "\author{Name \email{email} \orcid{orcid}}" gives a single footnote mark.
They can be part of \marg{name} or follow "\author"\marg{name} in the preamble, \eg "\author{Name}\email{email}\orcid{orcid}"; this also holds for "\thanks"\marg{text}, which gets its own footnote, and for editors and endorsers.
An "\email"\oarg{label}\marg{email} with label gets its own footnote, which can be referred to by the label.
After the title, "\email" and "\orcid" print the link in the text.
"\orcmail"\marg{orcid}\marg{email} is short for "\email"\marg{email}"\orcid"\marg{orcid}.
If editors are present, the author line receives a title that can be adjusted using "\authortitle", "\authortitlefont", "\preauthortitle", and "\postauthortitle".
The fonts can be adjusted using the "\authorfont"\marg{font} and "\affiliationfont"\marg{font} macros.
The following lines add \eg two authors with different affiliations
\begin{verbatim}
\author[affil1]{Author one \email{email one}}
\affiliation[affil1]{Affiliation one}
\author[affil1,affil2]{Author two \email{email two}}
\affiliation[affil2]{Affiliation two}
\end{verbatim}

\DescribeMacro{\editor}
\DescribeMacro{\endorser}
\DescribeMacro{\editorfont}
\DescribeMacro{\endorserfont}
\DescribeMacro{\preeditor}
\DescribeMacro{\posteditor}
\DescribeMacro{\preendorser}
\DescribeMacro{\postendorser}
Additionally the "\editor"\oarg{label}\marg{name} and "\endorser"\oarg{label}\marg{name} macros are provided that act similar to the "\author"\oarg{label}\marg{name} macros but place their content in a dedicated line.
Their font and surrounding material can be adjusted by "\editorfont"\marg{font}, "\endorserfont"\marg{font}, "\preeditor"\marg{code}, "\posteditor"\marg{code}, "\preendorser"\marg{code}, and "\postendorser"\marg{code}.
\DescribeMacro{\editortitle}
\DescribeMacro{\endorsertitle}
\DescribeMacro{\editortitlefont}
\DescribeMacro{\endorsertitlefont}
\DescribeMacro{\preeditortitle}
\DescribeMacro{\posteditortitle}
\DescribeMacro{\preendorsertitle}
\DescribeMacro{\postendorsertitle}
The titles of these lines can be adjusted using the "\editortitle"\marg{singular}\marg{plural} and "\endorsertitle"\marg{singular}\marg{plural} macros.
Their font and surrounding material can be adjusted with "\editortitlefont"\marg{font}, "\endorsertitlefont"\allowbreak\marg{font}, "\preeditortitle"\marg{code}, "\posteditortitle"\marg{code}, "\preendorsertitle"\marg{code}, and "\postendorsertitle"\marg{code}.

\subsection{Abstract}

\DescribeEnv{abstract}
\DescribeEnv{abstract*}
\DescribeMacro{\keywords}
The "abstract" environment is adjusted to not start with an indentation.
Journals expect the abstract in different places, and all of them are accepted: after "\maketitle", before it in the document, or in the preamble, as environment or in the command form "\abstract"\marg{text}.
An abstract given before the title is placed after it by "\maketitle".
The "abstract*" environment takes care of placing the title as well.
This is used in "twocolumn" documents if the abstract should span both columns.
The "\keywords"\marg{keywords} macro appends an italic keyword line to the abstract.

\DescribeMacro{\datefont}
The date is empty by default and can be formatted with "\datefont"\marg{font}.

\subsection{Preprint}

\DescribeMacro{\preprint}
\DescribeMacro{\arxiv}
\DescribeMacro{\placepreprint}
\DescribeMacro{\preprintfont}
The "\preprint"\marg{number} macro places a preprint number in the upper right corner of the title page.
The "\arxiv"\marg{id} macro adds an arXiv preprint link.
The "\placepreprint" hook places the preprint block on the first page, and "\preprintfont" changes its font.

\DescribeMacro{\reprint}
\DescribeMacro{\doi}
\DescribeMacro{\placereprint}
\DescribeMacro{\reprintfont}
The "\reprint"\marg{text} macro collects reprint or publication information for the upper left corner of the title page.
The "\doi"\marg{id} helper adds a DOI link.
The "\placereprint" hook places the reprint block on the first page, and "\reprintfont" changes its font.

% \ifshort
\printbibliography

\end{document}
%
%</documentation>
% \fi
%
% \StopEventually{
% \printbibliography
% \PrintChanges
% }
%
% \appendix
%
% \section{Implementation}
%
%<*package>
%
% \begin{macro}{\hep@url@text}
% Print a URL or email address in typewriter font without \software{hyperref}.
% As for "\nolinkurl", the text is printed verbatim, so that characters such as "_", "#", "&", and "~" need no escaping.
% The escaped characters "\%", "\#", "\&", "\_", "\$", and "\~" become the characters themselves.
% Macros in the text, such as the fields of \software{biblatex}, are expanded first.
% The \software{hep-text} and \software{hep-bibliography} packages define the same macro, so define it only once.
%    \begin{macrocode}
\ExplSyntaxOn
\cs_if_exist:NF \hep@url@text {
  \tl_new:N \l__hep_url_tl
  \cs_generate_variant:Nn \tl_replace_all:Nnn { NnV }
  \cs_new_protected:Npn \hep@url@text #1 {
    \tl_set:Nn \l__hep_url_tl {#1}
    \regex_replace_all:nnN { \cP. } { \cO\# } \l__hep_url_tl
    \protected@edef \l__hep_url_tl { \l__hep_url_tl }
    \tl_remove_all:Nn \l__hep_url_tl { \protect }
    \tl_replace_all:NnV \l__hep_url_tl { \# } \c_hash_str
    \tl_replace_all:NnV \l__hep_url_tl { \% } \c_percent_str
    \tl_replace_all:NnV \l__hep_url_tl { \& } \c_ampersand_str
    \tl_replace_all:NnV \l__hep_url_tl { \_ } \c_underscore_str
    \tl_replace_all:NnV \l__hep_url_tl { \$ } \c_dollar_str
    \tl_replace_all:NnV \l__hep_url_tl { \~ } \c_tilde_str
    \texttt { \tl_to_str:N \l__hep_url_tl }
  }
}
\ExplSyntaxOff
%    \end{macrocode}
% \end{macro}
%
% \begin{macro}{\online}
% \begin{macro}{\email}
% Provide the "\online"\marg{text}\marg{url} macro combining the features of the "\href" and the "\url" macros.
% Define a macro for typesetting emails.
%    \begin{macrocode}
\providecommand*{\online}[2]{\hep@url@text{#2}}%
\providecommand*{\hep@email}[1]{\online{mailto:#1}{#1}}
\providecommand*{\email}{\hep@email}
\AtBeginDocument{\@ifpackageloaded{hyperref}{%
    \renewcommand*{\online}[2]{\href{#1}{\nolinkurl{#2}}}%
  }{}
}
%    \end{macrocode}
% \end{macro}
% \end{macro}
%
% \begin{macro}{\hep@multi@ref}
% Transform reference list to list of references.
%    \begin{macrocode}
\newif\ifhep@first%
\newcommand*{\hep@multi@ref}[1]{%
  \hep@firsttrue%
  \forcsvlist{%
    \ifhep@first\hep@firstfalse\else\textsuperscript,\fi\ref%
  }{#1}%
}
%    \end{macrocode}
% \end{macro}
%
% \begin{macro}{\hep@contact}
% \begin{macro}{\hep@contact@add}
% \begin{macro}{\hep@contact@next}
% \begin{macro}{\hep@contact@check}
% \begin{macro}{\hep@contact@skip}
% \begin{macro}{\hep@contact@flush}
% Collect contact information, such as email addresses and ORCID links, for one footnote.
% After each item the next token is inspected; another "\email" or "\orcid" continues the footnote, everything else places it.
% Spaces are skipped while looking ahead and one space is restored after the footnote mark.
%    \begin{macrocode}
\newcommand*\hep@contact{}
\newif\ifhep@contact@space
\newcommand*\hep@contact@add[1]{%
  \ifdefempty{\hep@contact}{%
    \gdef\hep@contact{#1}%
  }{%
    \gappto\hep@contact{ #1}%
  }%
  \hep@contact@spacefalse%
  \hep@contact@next%
}
\newcommand*\hep@contact@next{%
  \futurelet\hep@contact@token\hep@contact@check%
}
\newcommand*\hep@contact@check{%
  \let\hep@contact@do\hep@contact@flush%
  \ifx\hep@contact@token\@sptoken%
    \let\hep@contact@do\hep@contact@skip%
  \fi%
  \ifx\hep@contact@token\email\let\hep@contact@do\relax\fi%
  \ifx\hep@contact@token\orcid\let\hep@contact@do\relax\fi%
  \hep@contact@do%
}
\@firstofone{\def\hep@contact@skip} {%
  \hep@contact@spacetrue\hep@contact@next%
}
\newcommand*\hep@contact@flush{%
  \unskip\expandafter\hep@contact@thanks\expandafter{\hep@contact}%
  \ifhep@contact@space\space\fi%
}
\newcommand*\hep@contact@thanks[1]{\gdef\hep@contact{}\thanks{#1}}
%    \end{macrocode}
% \end{macro}
% \end{macro}
% \end{macro}
% \end{macro}
% \end{macro}
% \end{macro}
%
% Redefine the email macro for the title page.
% In the preamble an email without label directly after "\author", "\editor", or "\endorser" is attached to this person.
% An email with label gets its own footnote.
%    \begin{macrocode}
\DeclareDocumentCommand{\email}{O{}m}{%
  \ifblank{#1}{%
    \hep@person@attach{\email{#2}}{\hep@contact@add{\hep@email{#2}}}%
  }{%
    \unskip\thanks[#1]{\hep@email{#2}}%
  }%
}
\providecommand*\hep@abstract@place{}
\AtBeginDocument{
  \let\hep@maketitle\maketitle
  \renewcommand*\maketitle{%
    \hep@maketitle\global\let\email\hep@email\global\let\orcid\orcidlink%
    \hep@abstract@place
  }
}
%    \end{macrocode}
%
% \subsection{Preprint and title}
%
% \subsubsection{Preprint and Reprint}
%
% \begin{macro}{\reprintfont}
% Define the "\reprintfont" macro.
%    \begin{macrocode}
\let\hep@reprint@font\relax
\newcommand*{\reprintfont}[1]{\def\hep@reprint@font{#1}}
\reprintfont{\scshape\small}
%    \end{macrocode}
% \end{macro}
%
% \begin{macro}{\reprint}
% \begin{macro}{\doi}
% Define the "\preprint" macro using the \software{varwidth} package \cite{varwidth}.
%    \begin{macrocode}
\let\hep@reprint\relax
\newcommand*\reprint[1]{%
  \@ifundefined{hep@reprint}{%
    \gdef\hep@reprint{#1}%
  }{%
    \apptocmd{\hep@reprint}{\\#1}{}{}%
  }%
}
\newcommand*\doi[1]{\reprint{\textsc{doi}: \online{https://doi.org/#1}{#1}}}
\newcommand*{\hep@reprint@box}{%
  \begin{varwidth}{\textwidth}%
    \hep@reprint@font\hep@reprint%
  \end{varwidth}%
}
%    \end{macrocode}
% \end{macro}
% \end{macro}
%
% \begin{macro}{\preprintfont}
% Define the "\preprintfont" macro.
%    \begin{macrocode}
\let\hep@preprint@font\relax
\newcommand*{\preprintfont}[1]{\def\hep@preprint@font{#1}}
\preprintfont{\scshape\small\raggedleft}
%    \end{macrocode}
% \end{macro}
%
% \begin{macro}{\preprint}
% \begin{macro}{\arxiv}
% Define the "\preprint" macro using the \software{varwidth} package \cite{varwidth}.
%    \begin{macrocode}
\RequirePackage{etoolbox}
\let\hep@preprint\relax
\newcommand*\preprint[1]{%
  \@ifundefined{hep@preprint}{%
    \gdef\hep@preprint{#1}%
  }{%
    \apptocmd{\hep@preprint}{\\#1}{}{}%
  }%
}
\newcommand*\arxiv[1]{%
  \preprint{\textsc{arXiv}: \online{https://arxiv.org/abs/#1}{#1}}%
}
\RequirePackage{varwidth}
\newcommand*{\hep@preprint@box}{%
  \begin{varwidth}{\textwidth}%
    \hep@preprint@font\hep@preprint%
  \end{varwidth}%
}
%    \end{macrocode}
% \end{macro}
% \end{macro}
%
% \begin{macro}{\placepreprint}
% \begin{macro}{\placereprint}
% Places a preprint number in the top right corner of the title page using the \software{calc} \cite{calc} and \software{atbegshi} \cite{atbegshi} packages.
% The box dimensions are measured into lengths before they enter the "\put" coordinates, because \software{calc} expressions such as "\heightof" are no longer accepted there.
%    \begin{macrocode}
\RequirePackage{calc}
\RequirePackage{atbegshi}
\newlength{\hep@box@width}
\newlength{\hep@box@height}
\newcommand*{\placepreprint}{%
  \AtBeginShipoutFirst{%
    \settowidth{\hep@box@width}{\normalfont\hep@preprint@box}%
    \settoheight{\hep@box@height}{\normalfont\hep@preprint@box}%
    \put(
      \textwidth+\oddsidemargin-\hep@box@width,
      -2pt-\topmargin-\hep@box@height
    ){\normalfont\hep@preprint@box}
  }
}
\newcommand*{\placereprint}{%
  \AtBeginShipoutFirst{%
    \settoheight{\hep@box@height}{\normalfont\hep@reprint@box}%
    \put(
      \oddsidemargin      ,
      -2pt-\topmargin-\hep@box@height
    ){\normalfont\hep@reprint@box}
  }
}
%    \end{macrocode}
% \end{macro}
% \end{macro}
%
% \subsubsection{Series}
%
% \begin{macro}{\series}
% Define a series title.
%    \begin{macrocode}
\newcommand{\preseries}[1]{\def\hep@pre@series{#1}}
\newcommand*{\series}[1]{\def\hep@series{#1}}
\newcommand{\postseries}[1]{\def\hep@post@series{#1}}
%    \end{macrocode}
% \end{macro}
%
% \begin{macro}{\seriesfont}
% Allow to change the fontface of the series title.
%    \begin{macrocode}
\let\hep@series@font\relax
\newcommand*{\seriesfont}[1]{\def\hep@series@font{#1}}
%    \end{macrocode}
% \end{macro}
%
% Set the default series title layout.
%    \begin{macrocode}
\preseries{\begin{center}\Large\hep@series@font}
\postseries{\par\end{center}}
%    \end{macrocode}
%
% \subsubsection{Title}
%
% Extend the title using the \software{titling} package \cite{titling}.
% Fix the footnote indent.
%    \begin{macrocode}
\RequirePackage{titling}
\setlength{\thanksmarkwidth}{1.5em}
%    \end{macrocode}
%
% \begin{macro}{\maketitlehooka}
% Place the preprint and the series title using "\maketitlehooka".
%    \begin{macrocode}
\renewcommand*{\maketitlehooka}{%
  \placereprint\placepreprint\vspace{-\bigskipamount}%
  \@ifundefined{hep@series}{}{%
    \hep@pre@series\hep@series\hep@post@series%
  }%
%   \vspace{-\bigskipamount}%
}
%    \end{macrocode}
% \end{macro}
%
% \begin{macro}{\titlefont}
% Allow to change the fontface of the the title.
%    \begin{macrocode}
\let\hep@title@font\relax
\newcommand*{\titlefont}[1]{\def\hep@title@font{#1}}
%    \end{macrocode}
% \end{macro}
% Set default title layout.
%    \begin{macrocode}
\pretitle{\begin{center}\LARGE\hep@title@font}
\posttitle{\par\end{center}}
%    \end{macrocode}

% \subsubsection{Subtitle}
%
% \begin{macro}{\subtitle}
% Define a subtitle.
%    \begin{macrocode}
\newcommand{\presubtitle}[1]{\def\hep@pre@sub@title{#1}}
\newcommand*{\subtitle}[1]{\def\hep@sub@title{#1}}
\newcommand{\postsubtitle}[1]{\def\hep@post@sub@title{#1}}
%    \end{macrocode}
% \end{macro}

% \begin{macro}{\subtitlefont}
% Allow to change the fontface of the subtitle.
%    \begin{macrocode}
\let\hep@subtitle@font\relax
\newcommand*{\subtitlefont}[1]{\def\hep@subtitle@font{#1}}
%    \end{macrocode}
% \end{macro}
%
% Set default subtitle layout.
%    \begin{macrocode}
\presubtitle{\begin{center}\Large\hep@subtitle@font}
\postsubtitle{\par\end{center}}
%    \end{macrocode}
%
% \subsection{Authors, Editors and Endorsers}
%
% Enable the handling of multiple authors with different affiliations using the \software{authblk} package \cite{authblk}.
% \software{authblk} defines the state of its author blocks only when an author starts a new block, which is not the case for the first author if an editor or endorser comes first.
%    \begin{macrocode}
\RequirePackage{authblk}
\providecommand*\AB@as{}
\providecommand*\AB@las{}
\providecommand*\AB@lasx{\protect\Authand}
\providecommand*\AB@au@str{}
%    \end{macrocode}
%
% \subsubsection{Deferred entries}
%
% \begin{macro}{\hep\_person:nnn}
% \begin{macro}{\hep\_person\_commit:}
% \begin{macro}{\hep\_person\_attach:nn}
% In the preamble an author, editor, or endorser is only stored, and added to its list by the next "\author", "\editor", "\endorser", or "\affiliation", or at the beginning of the document.
% Until then, "\email", "\orcid", "\orcmail", and "\thanks" without label are appended to the stored name, so that they can follow "\author"\marg{name} instead of being part of \marg{name}.
% In the document the entries are added immediately.
%    \begin{macrocode}
\ExplSyntaxOn
\tl_new:N \g__hep_person_tl
\tl_new:N \g__hep_person_label_tl
\tl_new:N \g__hep_person_name_tl
\cs_new_protected:Npn \__hep_person_add:nnn #1#2#3
  {
    \str_if_eq:nnTF {#1} { author }
      {
        \hep_add_to_authors:n {#3}
        \hep_author [#2] {#3}
      }
      { \hep@person@add {#1} {#2} {#3} }
  }
\cs_generate_variant:Nn \__hep_person_add:nnn { VVV }
\cs_new_protected:Npn \hep_person_commit:
  {
    \tl_if_empty:NF \g__hep_person_tl
      {
        \__hep_person_add:VVV \g__hep_person_tl
          \g__hep_person_label_tl \g__hep_person_name_tl
        \tl_gclear:N \g__hep_person_tl
      }
  }
\cs_new_protected:Npn \hep_person:nnn #1#2#3
  {
    \hep_person_commit:
    \tl_gset:Nn \g__hep_person_tl {#1}
    \tl_gset:Nn \g__hep_person_label_tl {#2}
    \tl_gset:Nn \g__hep_person_name_tl {#3}
    \cs_if_eq:NNT \@nodocument \relax { \hep_person_commit: }
  }
\cs_new_protected:Npn \hep_person_attach:nn #1#2
  {
    \tl_if_empty:NTF \g__hep_person_tl {#2}
      { \tl_gput_right:Nn \g__hep_person_name_tl {#1} }
  }
\cs_new_eq:NN \hep@person \hep_person:nnn
\cs_new_eq:NN \hep@person@commit \hep_person_commit:
\cs_new_eq:NN \hep@person@attach \hep_person_attach:nn
\AddToHook{begindocument/before}{\hep_person_commit:}
\ExplSyntaxOff
\RequirePackage{xpatch}
\xpretocmd{\affil}{\hep@person@commit}{}{}
%    \end{macrocode}
% \end{macro}
% \end{macro}
% \end{macro}
%
% \subsubsection{Editors and endorsers}
%
% \begin{macro}{\hep@person@add}
% Add a name with an optional list of labels to the list of editors or endorsers, similar to the "\author" macro of \software{authblk}.
% The first argument is the kind of person.
%    \begin{macrocode}
\newcommand*\hep@person@add[3]{%
  \ifnewaffil%
    \addtocounter{affil}{1}%
    \edef\AB@thenote{\arabic{affil}}%
  \fi%
  \if\relax#2\relax\else\setcounter{Maxaffil}{0}\fi%
  \ifnum\value{#1s}>1\relax%
    \@namedef{hep@#1@sep\number\value{#1s}}{\Authsep}%
  \fi%
  \addtocounter{#1s}{1}%
  \begingroup%
    \let\protect\@unexpandable@protect%
    \let\and\AB@pand%
    \def\thanks{\protect\thanks}\def\footnote{\protect\footnote}%
    \let\\\authorcr%
    \@temptokena=\expandafter\expandafter\expandafter{%
      \csname hep@#1list\endcsname%
    }%
    \expandafter\xdef\csname hep@#1list\endcsname{%
      \the\@temptokena%
      \protect\@nameuse{hep@#1@sep\number\value{#1s}}%
      \protect\Authfont\hep@author@name{#3}%
      \if\relax#2\relax\else\protect\hep@multi@ref{#2}\fi%
    }%
  \endgroup%
  \ifnum\value{#1s}>2\relax%
    \@namedef{hep@#1@sep\number\value{#1s}}{\Authands}%
  \fi%
  \newaffilfalse%
}
%    \end{macrocode}
% \end{macro}
%
% \begin{macro}{\hep@print@person}
% Print the line of the given kind of persons, if there are any.
%    \begin{macrocode}
\newcommand*\hep@print@person[1]{%
  \expandafter\ifx\csname hep@#1list\endcsname\AB@empty\else%
    \@nameuse{hep@pre@#1}\@nameuse{hep@#1list}\@nameuse{hep@post@#1}%
  \fi%
}
%    \end{macrocode}
% \end{macro}
%
% \begin{macro}{\hep@new@setter}
% Define a macro \marg{name} that stores its argument in the macro \marg{internal}.
% The argument may contain paragraphs.
% "\def" is used, since "\newcommand" refuses names starting with "end", such as "\endorserfont".
%    \begin{macrocode}
\newcommand*\hep@new@setter[2]{%
  \expandafter\long\expandafter\def\csname#1\endcsname##1{%
    \@namedef{#2}{##1}%
  }%
}
%    \end{macrocode}
% \end{macro}
%
% \begin{macro}{\hep@new@person}
% Define the user interface for a kind of person \marg{kind} with the line titles \marg{singular} and \marg{plural}:
% the counter \meta{kind}"s", the macro "\"\meta{kind}\oarg{label}\marg{name}, the font macros "\"\meta{kind}"font" and "\"\meta{kind}"titlefont", the title macro "\"\meta{kind}"title"\marg{singular}\marg{plural}, and the hooks "\pre"\meta{kind}, "\post"\meta{kind}, "\pre"\meta{kind}"title", and "\post"\meta{kind}"title".
% By default the line is centred and starts with the italic title, followed by a colon.
%    \begin{macrocode}
\newcommand*\hep@new@person[3]{%
  \newcounter{#1s}%
  \expandafter\let\csname hep@#1list\endcsname\AB@empty%
  \@namedef{hep@#1@sep1}{}%
  \@namedef{hep@#1@sep2}{\Authand}%
  \expandafter\NewDocumentCommand\csname#1\endcsname{O{}m}{%
    \hep@person{#1}{##1}{##2}%
  }%
  \hep@new@setter{#1font}{hep@#1@font}%
  \hep@new@setter{pre#1}{hep@pre@#1}%
  \hep@new@setter{post#1}{hep@post@#1}%
  \hep@new@setter{#1titlefont}{hep@#1@title@font}%
  \hep@new@setter{pre#1title}{hep@pre@#1@title}%
  \hep@new@setter{post#1title}{hep@post@#1@title}%
  \expandafter\def\csname#1title\endcsname##1##2{%
    \@namedef{hep@#1@title}{##1}%
    \@namedef{hep@#1@title@pl}{##2}%
  }%
  \@namedef{hep@#1@font}{}%
  \@nameuse{#1title}{#2}{#3}%
  \@nameuse{#1titlefont}{\itshape}%
  \@nameuse{pre#1title}{\@nameuse{hep@#1@title@font}}%
  \@nameuse{post#1title}{: }%
  \@nameuse{pre#1}{%
    \begin{center}%
      \large\@nameuse{hep@#1@font}\lineskip.5em%
      \begin{tabular}[t]{c}{%
        \@nameuse{hep@pre@#1@title}%
        \ifnum\value{#1s}>1\relax%
          \@nameuse{hep@#1@title@pl}%
        \else%
          \@nameuse{hep@#1@title}%
        \fi%
        \@nameuse{hep@post@#1@title}%
      }%
  }%
  \@nameuse{post#1}{\end{tabular}\par\end{center}}%
}
%    \end{macrocode}
% \end{macro}
%
% \begin{macro}{\editor}
% \begin{macro}{\editorfont}
% \begin{macro}{\preeditor}
% \begin{macro}{\posteditor}
% \begin{macro}{\editortitle}
% \begin{macro}{\editortitlefont}
% \begin{macro}{\preeditortitle}
% \begin{macro}{\posteditortitle}
% \begin{macro}{\endorser}
% \begin{macro}{\endorserfont}
% \begin{macro}{\preendorser}
% \begin{macro}{\postendorser}
% \begin{macro}{\endorsertitle}
% \begin{macro}{\endorsertitlefont}
% \begin{macro}{\preendorsertitle}
% \begin{macro}{\postendorsertitle}
% Define editors and endorsers.
%    \begin{macrocode}
\hep@new@person{editor}{Editor}{Editors}
\hep@new@person{endorser}{Endorser}{Endorsers}
%    \end{macrocode}
% \end{macro}
% \end{macro}
% \end{macro}
% \end{macro}
% \end{macro}
% \end{macro}
% \end{macro}
% \end{macro}
% \end{macro}
% \end{macro}
% \end{macro}
% \end{macro}
% \end{macro}
% \end{macro}
% \end{macro}
% \end{macro}
%
% \begin{macro}{\maketitlehookb}
% Show subtitle and editor.
%    \begin{macrocode}
\renewcommand*{\maketitlehookb}{%
  \@ifundefined{hep@sub@title}{}{%
    \hep@pre@sub@title\hep@sub@title\hep@post@sub@title%
  }%
  \smallskip%
  \hep@print@person{editor}%
}
%    \end{macrocode}
% \end{macro}
%
% \subsubsection{Authors}
%
% \begin{macro}{\author}
% Allow absent author field.
%    \begin{macrocode}
% \author{}
%    \end{macrocode}
% \end{macro}
%
% \begin{macro}{\hep@author@name}
% With many affiliations \software{authblk} typesets the authors as a centred paragraph, in which every break costs the same.
% Each name is boxed, so that the lines break only between the names, at the commas or around \enquote{and}, and names are neither broken nor hyphenated.
% Inside "\maketitle" "\thanks", "\email", and "\orcid" only place a mark in the box and store the footnote text; "\footnote" is made to do the same, since its text would be lost in the box.
% The final "\unskip" removes a space after the name, such as the one restored after a contact mark.
%    \begin{macrocode}
\DeclareRobustCommand*\hep@author@name[1]{%
  \mbox{\let\footnote\thanks#1\unskip}%
}
%    \end{macrocode}
% \end{macro}
%
% Switch authblk to a label ref system for affiliations using the \software{xpatch} package \cite{xpatch}.
%    \begin{macrocode}
\RequirePackage{xpatch}
\xpatchcmd{\author}{%
  \protect\Authfont#2\AB@authnote{\AB@note}%
}{%
  \protect\Authfont\hep@author@name{#2}%
  \if\relax#1\relax\else\unskip\protect\hep@multi@ref{\AB@note}\fi%
}{}{}
%    \end{macrocode}
%
% \begin{macro}{\authorfont}
% Allow to change the fontface of the individual parts of the title.
%    \begin{macrocode}
\let\hep@author@font\relax
\newcommand*{\authorfont}[1]{\def\hep@author@font{#1}}
\renewcommand*\Authfont{\hep@author@font}
%    \end{macrocode}
% \end{macro}
%
% Set default author fontface.
%    \begin{macrocode}
\newcommand*{\authortitle}[2]{%
  \def\hep@author@title{#1}%
  \def\hep@author@title@pl{#2}%
}
\newcommand*{\authortitlefont}[1]{\def\hep@author@title@font{#1}}
\newcommand{\preauthortitle}[1]{\def\hep@pre@author@title{#1}}
\newcommand{\postauthortitle}[1]{\def\hep@post@author@title{#1}}
\authortitle{Author}{Authors}
\authortitlefont{\itshape}
\preauthortitle{\hep@author@title@font}
\postauthortitle{: }
\preauthor{%
  \begin{center}%
    \large\hep@author@font\lineskip.5em%
    \begin{tabular}[t]{c}{%
      \ifnum\value{editors}>0\relax%
        \hep@pre@author@title%
        \ifnum\value{authors}>1\relax%
          \hep@author@title@pl\else\hep@author@title%
        \fi\hep@post@author@title%
      \fi%
    }%
}
\postauthor{\end{tabular}\par\end{center}}
%    \end{macrocode}
%
% \subsubsection{Affiliation}
%
% Patch the "\affiliation" macro to comply with the label ref system.
%    \begin{macrocode}
\newcounter{affiliation}
\renewcommand*{\theaffiliation}{%
  \textsuperscript{\normalfont\alph{affiliation}}%
}
\xpatchcmd{\affil}{%
  \AB@affilnote{\AB@note}%
}{%
  \protect\refstepcounter{affiliation}\protect\label{\AB@note}%
  \if\relax#1\relax\else\protect\theaffiliation\fi%
}{}{}
%    \end{macrocode}
%
% \begin{macro}{\affiliationfont}
% Allow to change the fontface of affiliation.
%    \begin{macrocode}
\let\hep@affiliation@font\relax
\newcommand*{\affiliationfont}[1]{\def\hep@affiliation@font{#1}}
% \newif\ifhep@lining\hep@liningtrue
% \ifhep@lining
  \renewcommand*{\Affilfont}{\small\hep@affiliation@font}
% \else
%   \renewcommand*{\Affilfont}{\small\ostyle\hep@affiliation@font}
% \fi
%    \end{macrocode}
% \end{macro}
%
% \begin{macro}{\affiliation}
% Define the "\affiliation" macro, ensure that linebreaks happen after a comma.
% The penalty is robust, so that \software{authblk} stores it unexpanded and the column mode is checked when the affiliation is typeset.
% With the flexible margins of the affiliation block a break with a penalty of $-50$ has negative demerits, so on its own the comma penalty would add lines.
% Therefore each affiliation, including its label, is measured twice in the current paragraph shape: once without comma penalties and with a high "\linepenalty", which gives the smallest number of lines, and once as it is typeset.
% "\looseness" then asks for the smallest number of lines, among which the breaks after commas are preferred.
% Footnotes are dropped while measuring, so that they are not counted twice, and "\everypar" is cleared and the paragraph ends with the primitive "\par", so that the pending list item of the "center" environment is left for the affiliation itself.
%    \begin{macrocode}
\newcommand*\hep@penalty{\if@twocolumn85\else50\fi}
\newif\ifhep@affil@break\hep@affil@breaktrue
\DeclareRobustCommand*\hep@comma@penalty{%
  \ifhep@affil@break\penalty-\hep@penalty\relax\fi%
}
\newcommand*\hep@active@comma{,\hep@comma@penalty}
\newcommand*\hep@cat@comma@active{\catcode`\,\active}
{\hep@cat@comma@active\gdef,{\hep@active@comma}}
\newcount\hep@affil@lines
\newcommand\hep@affil@measure[2]{%
  \setbox\z@\vbox{%
    \let\footnote\@gobble\let\thanks\@gobble%
    \let\footnotemark\@empty%
    \everypar{}\looseness\z@\noindent%
    \if\relax#1\relax\else\theaffiliation\fi#2\@@par%
    \global\hep@affil@lines\prevgraf%
  }%
}
\DeclareRobustCommand\hep@affil@text[2]{%
  \begingroup%
    \hep@affil@measure{#1}{#2}%
    \count@\hep@affil@lines%
    \hep@affil@breakfalse\linepenalty\@m%
    \hep@affil@measure{#1}{#2}%
    \advance\count@-\hep@affil@lines%
    \count@-\count@%
  \expandafter\endgroup\expandafter\looseness\the\count@\relax%
  #2%
}
\newcommand*\hep@affil[1]{\endgroup\affil{\hep@affil@text{}{#1}}}
\def\hep@affil@opt[#1]#2{%
  \endgroup\affil[#1]{\hep@affil@text{#1}{#2}}%
}
\DeclareRobustCommand\hep@affiliation{%
  \@ifnextchar[{\hep@affil@opt}{\hep@affil}%
}
\newcommand*{\affiliation}{%
  \begingroup\hep@cat@comma@active\hep@affiliation%
}
%    \end{macrocode}
% \end{macro}
%
% Place endorser and affiliation.
% The affiliations are typeset with flexible margins, so that the line breaks after commas are preferred; this setting is local to the affiliation block.
%    \begin{macrocode}
\renewcommand*{\maketitlehookc}{%
  \hep@print@person{endorser}%
%       \\[\affilsep]
  \ifx\AB@affillist\AB@empty\else%
    \ifnum\value{affil}>\value{Maxaffil}\relax%
      \begingroup%
        \@flushglue=0pt plus .5\linewidth%
        \begin{center}\AB@affillist\end{center}%
      \endgroup%
    \fi%
  \fi%
}
%    \end{macrocode}
%
% \begin{macro}{\@author}
% Ensure that affiliation is not set directly below author.
% Without affiliations the list of all authors is used, since \software{authblk} keeps only the last one in this case; it falls back to the author given before \software{hep-title} has been loaded.
% The "\rlap" redefinition lives in a helper macro, so that "\@author" contains no parameter characters; \software{hyperref} copies "\@author" with the "pdfusetitle" option.
%    \begin{macrocode}
\newcommand*{\hep@no@rlap}{\def\rlap##1{##1}}
\def\@author{%
  \ifx\AB@affillist\AB@empty%
    \ifx\AB@authlist\AB@empty\AB@author\else\AB@authlist\fi%
  \else%
    \ifnum\value{affil}>\value{Maxaffil}\relax%
      \hep@no@rlap%
      \AB@authlist%
%       \\[\affilsep]\AB@affillist%
    \else%
      \AB@authors%
    \fi%
  \fi%
}
%    \end{macrocode}
% \end{macro}
% \subsubsection{PDF strings}
%
% \begin{macro}{\_\_hep\_pdf\_gobble\_opt:w}
% \begin{macro}{\_\_hep\_pdf\_drop\_opt:w}
% \begin{macro}{\_\_hep\_pdf\_newline:w}
% Expandable replacements for PDF strings.
% The first one removes a footnote-like command with an optional label, the second one keeps the mandatory argument and the third one turns a line break with optional star and length into a space.
%    \begin{macrocode}
\ExplSyntaxOn
\cs_new:Npn \__hep_pdf_gobble_opt:w #1
  { \str_if_eq:nnT {#1} { [ } { \__hep_pdf_gobble_opt_aux:w } }
\cs_new:Npn \__hep_pdf_gobble_opt_aux:w #1 ] #2 { }
\cs_new:Npn \__hep_pdf_drop_opt:w #1
  { \str_if_eq:nnTF {#1} { [ } { \__hep_pdf_drop_opt_aux:w } {#1} }
\cs_new:Npn \__hep_pdf_drop_opt_aux:w #1 ] #2 {#2}
\cs_new:Npn \__hep_pdf_newline:w #1
  {
    \str_if_eq:nnTF {#1} { * }
      { \__hep_pdf_newline:w }
      {
        \str_if_eq:nnTF {#1} { [ }
          { \__hep_pdf_newline_aux:w } { ~ #1 }
      }
  }
\cs_new:Npn \__hep_pdf_newline_aux:w #1 ] #2 { ~ #2 }
%    \end{macrocode}
% \end{macro}
% \end{macro}
% \end{macro}
%
% \subsubsection{PDF document information}
%
% \begin{macro}{\g\_hep\_authors\_seq}
% \begin{macro}{\hep\_author}
% Collect the authors for the PDF document information.
%    \begin{macrocode}
\seq_new:N \g_hep_authors_seq
\cs_new_eq:NN \hep_author \author
\tl_new:N \l__hep_pdf_tl
%    \end{macrocode}
% \end{macro}
% \end{macro}
%
% \begin{macro}{\_\_hep\_pdf\_normalise:N}
% Remove the spaces in front of the commands that are dropped in PDF strings and around line breaks and "\and", and trim the spaces at both ends.
%    \begin{macrocode}
\cs_new_protected:Npn \__hep_pdf_normalise:N #1
  {
    \clist_map_inline:nn
      { \email, \thanks, \footnote, \orcid, \orcidlink, \orcmail, \\, \and }
      { \tl_replace_all:Nnn #1 { ~ ##1 } {##1} }
    \tl_replace_all:Nnn #1 { \\ ~ } { \\ }
    \tl_replace_all:Nnn #1 { \and ~ } { \and }
    \tl_trim_spaces:N #1
  }
%    \end{macrocode}
% \end{macro}
%
% \begin{macro}{\hep\_add\_to\_authors:n}
% Store the author unexpanded, it is converted into a PDF string by \software{hyperref} when the title is made.
%    \begin{macrocode}
\cs_new_protected:Npn \hep_add_to_authors:n #1
  {
    \tl_set:Nn \l__hep_pdf_tl {#1}
    \__hep_pdf_normalise:N \l__hep_pdf_tl
    \seq_gput_right:NV \g_hep_authors_seq \l__hep_pdf_tl
  }
%    \end{macrocode}
% \end{macro}
%
% \begin{macro}{\author}
% Patch author macro, the author is added to the lists by \cs{hep\_person\_commit:}.
%    \begin{macrocode}
\RenewDocumentCommand{\author}{O{}m}{
  \hep_person:nnn {author} {#1} {#2}
}
%    \end{macrocode}
% \end{macro}
%
% \begin{macro}{\hep\_authors:}
% \begin{macro}{\hep@authors}
% Expandable list of the authors.
%    \begin{macrocode}
\cs_new:Npn \hep_authors:
  {
    \seq_use:Nnnn \g_hep_authors_seq
      { ~ and ~ } { , ~ } { , ~ and ~ }
  }
\cs_new_eq:NN \hep@authors \hep_authors:
%    \end{macrocode}
% \end{macro}
% \end{macro}
%
% \begin{macro}{\_\_hep\_pdf\_info:nn}
% Set the PDF document information entry \marg{key} to \marg{value} unless it has been set by the user, \ie unless it is neither empty nor the value set by a previous call.
%    \begin{macrocode}
\tl_new:N \g__hep_pdf_author_tl
\tl_new:N \g__hep_pdf_title_tl
\tl_new:N \g__hep_pdf_keywords_tl
\cs_new_protected:Npn \__hep_pdf_info:nn #1#2
  {
    \bool_lazy_or:nnT
      { \tl_if_empty_p:c { @pdf #1 } }
      { \tl_if_eq_p:cc { @pdf #1 } { g__hep_pdf_ #1 _tl } }
      {
        \hypersetup { pdf #1 = {#2} }
        \tl_gset_eq:cc { g__hep_pdf_ #1 _tl } { @pdf #1 }
      }
  }
\cs_generate_variant:Nn \__hep_pdf_info:nn { nV }
%    \end{macrocode}
% \end{macro}
%
% \begin{macro}{\_\_hep\_pdf\_info:}
% Set the PDF title, author and keywords from "\title", "\subtitle", "\author" and "\keywords" similar to the "pdfusetitle" option of \software{hyperref}.
% Footnote-like commands are dropped, line breaks become spaces and "\and" becomes a comma.
%    \begin{macrocode}
\cs_new_protected:Npn \__hep_pdf_info:
  {
    \cs_gset_eq:NN \__hep_pdf_saved_hook: \pdfstringdefPreHook
    \pdfstringdefDisableCommands
      {
        \cs_set_eq:NN \email \__hep_pdf_gobble_opt:w
        \cs_set_eq:NN \thanks \__hep_pdf_gobble_opt:w
        \cs_set_eq:NN \footnote \__hep_pdf_gobble_opt:w
        \cs_set_eq:NN \orcidlink \use_none:n
        \cs_set_eq:NN \orcid \use_none:n
        \cs_set_eq:NN \orcmail \use_none:nn
        \cs_set_eq:NN \\ \__hep_pdf_newline:w
        \cs_set_eq:NN \newline \__hep_pdf_newline:w
        \cs_set:Npn \and { , ~ }
      }
    \tl_if_in:NnF \@title { \@latex@error }
      {
        \tl_set_eq:NN \l__hep_pdf_tl \@title
        \__hep_pdf_normalise:N \l__hep_pdf_tl
        \@ifundefined { hep@sub@title } { }
          { \tl_put_right:Nn \l__hep_pdf_tl { : ~ \hep@sub@title } }
        \tl_put_right:Nn \l__hep_pdf_tl { \@empty }
        \__hep_pdf_info:nV { title } \l__hep_pdf_tl
      }
    \seq_if_empty:NF \g_hep_authors_seq
      {
        \tl_set:Ne \l__hep_pdf_tl { \hep_authors: \exp_not:N \@empty }
        \__hep_pdf_info:nV { author } \l__hep_pdf_tl
      }
    \@ifundefined { hep@key@words } { }
      { \__hep_pdf_info:nn { keywords } { \hep@key@words } }
    \cs_gset_eq:NN \pdfstringdefPreHook \__hep_pdf_saved_hook:
  }
\AddToHook{cmd/maketitle/before}[hep-title/pdfinfo]{
  \@ifpackageloaded{hyperref}{\__hep_pdf_info:}{}
}
\ExplSyntaxOff
%    \end{macrocode}
% \end{macro}
%
% \begin{macro}{\orcid}
% \begin{macro}{\orcmail}
% Load the \software{orcidlink} \cite{orcidlink} package, which also loads \software{hyperref}, and make the logo monochrome.
% It is loaded after "\author" has been redefined, so that the "pdfusetitle" option of \software{hyperref} still sees the authors.
% Define the "\orcid" macro that places the ORCID link in the footnote of the contact information, like "\email".
% The "\orcmail" macro is kept for backward compatibility.
%    \begin{macrocode}
\RequirePackage{orcidlink}
\colorlet{orcidlogocol}{black}
\tikzset{
  orcidlogo/.pic={
    \fill[orcidlogocol] svg{M256,128c0,70.7-57.3,128-128,128C57.3,256,0,198.7,0,128C0,57.3,57.3,0,128,0C198.7,0,256,57.3,256,128z};
    \fill[white] svg{M236,128 c0,59.6-48.4,108-108,108 C68.4,236,20,187.6,20,128 C20,68.4,68.4,20,128,20 C187.6,20,236,68.4,236,128z};
    \fill[black] svg{M86.3,186.2H70.9V79.1h15.4v48.4V186.2z}
    svg{M108.9,79.1h41.6c39.6,0,57,28.3,57,53.6c0,27.5-21.5,53.6-56.8,53.6h-41.8V79.1z M124.3,172.4h24.5c34.9,0,42.9-26.5,42.9-39.7c0-21.5-13.7-39.7-43.7-39.7h-23.7V172.4z}
    svg{M88.7,56.8c0,5.5-4.5,10.1-10.1,10.1c-5.6,0-10.1-4.6-10.1-10.1c0-5.6,4.5-10.1,10.1-10.1C84.2,46.7,88.7,51.3,88.7,56.8z};
  }
}
\DeclareDocumentCommand{\orcid}{m}{%
  \hep@person@attach{\orcid{#1}}{\hep@contact@add{\orcidlinkc{#1}}}%
}
\NewDocumentCommand{\orcmail}{mm}{\email{#2}\orcid{#1}}
%    \end{macrocode}
% \end{macro}
% \end{macro}
%
% In PDF strings such as bookmarks the "\email" macro shows the address and the "\orcid" and "\orcmail" macros are dropped.
%    \begin{macrocode}
\ExplSyntaxOn
\pdfstringdefDisableCommands{
  \cs_set_eq:NN \email \__hep_pdf_drop_opt:w
  \cs_set_eq:NN \orcid \use_none:n
  \cs_set_eq:NN \orcmail \use_none:nn
  \cs_set_eq:NN \hep@author@name \use:n
}
\ExplSyntaxOff
%    \end{macrocode}
%
% \subsection{Date and Abstract}
%
% \subsubsection{Date}
%
% \begin{macro}{\date}
% Allow absent date field.
%    \begin{macrocode}
\date{\vspace{-4ex}}
%    \end{macrocode}
% \end{macro}
%
% \begin{macro}{\datefont}
% Allow to change the fontface of the individual parts of the title.
%    \begin{macrocode}
\let\hep@date@font\relax
\newcommand*{\datefont}[1]{\def\hep@date@font{#1}}
%    \end{macrocode}
% \end{macro}
%
% Set the default "date" fontface.
%    \begin{macrocode}
\predate{\begin{center}\hep@date@font}
\postdate{\par\end{center}}
%    \end{macrocode}
%
% \subsubsection{Abstract}
%
% \begin{environment}{abstract*}
% Add a "abstract*" environment for two column mode taking also care of placing the title using the \software{environ} \cite{environ} and \software{abstract} \cite{abstract} packages.
% It is defined first, since the \software{abstract} package redefines the "abstract" environment.
%    \begin{macrocode}
\@ifundefined{abstract}{}{%
  \if@twocolumn
    \RequirePackage{environ}
    \RequirePackage{abstract}
    \renewcommand*{\abstitleskip}{-3ex}
    \NewEnviron{abstract*}{%
      \twocolumn[\maketitle\vspace{-5ex}%
      \begin{onecolabstract}\noindent\BODY\end{onecolabstract}%
      \vspace{.5cm}]\saythanks%
    }%
  \else
    \newenvironment{abstract*}{%
      \maketitle\begin{abstract}%
    }{%
      \end{abstract}%
    }
  \fi
%    \end{macrocode}
% \end{environment}
%
% \begin{environment}{abstract}
% Adjust the "abstract" environment to not start with indentation.
%    \begin{macrocode}
  \let\hep@abstract\abstract%
  \newcommand*\hep@abstract@begin{\hep@abstract\noindent\ignorespaces}%
  \let\hep@abstract@end\endabstract
  \let\hep@key@words\relax
  \newcommand*{\keywords}[1]{\def\hep@key@words{#1}}
  \AtBeginDocument{\@ifundefined{hep@key@words}{}{%
    \let\hep@end@abstract\hep@abstract@end%
    \renewcommand*\hep@abstract@end{
      \par\medskip\noindent\textit{Keywords: \hep@key@words}%
      \hep@end@abstract%
    }%
  }}
%    \end{macrocode}
%
% \begin{macro}{\ifhep@title@made}
% \begin{macro}{\hep@abstract@store}
% \begin{macro}{\hep@abstract@collect}
% Journals expect the abstract in different places: after "\maketitle" (\software{article}, \software{svjour3}), before it in the document (REV\hologo{TeX}, \software{amsart}), or in the preamble, either as environment or in the command form "\abstract"\marg{text} (\software{JHEP}, \software{JCAP}, PoS).
% Accept all of them: before the title is made the abstract is only stored, collected by the \software{environ} package \cite{environ} if it is an environment, and "\maketitle" places it after the title.
%    \begin{macrocode}
  \newif\ifhep@title@made
  \newif\ifhep@abstract@stored
  \RequirePackage{environ}
  \newcommand*\hep@abstract@name{abstract}
  \newcommand\hep@abstract@store[1]{%
    \gdef\hep@abstract@text{#1}%
    \global\hep@abstract@storedtrue
  }
  \newcommand*\hep@abstract@collect{%
    \Collect@Body\hep@abstract@store
  }
  \newif\ifhep@abstract@collecting
  \renewcommand*\abstract{%
    \ifhep@title@made
      \hep@abstract@collectingfalse
      \let\hep@abstract@next\hep@abstract@begin
    \else
      \ifx\@currenvir\hep@abstract@name
        \hep@abstract@collectingtrue
        \let\hep@abstract@next\hep@abstract@collect
      \else
        \let\hep@abstract@next\hep@abstract@store
      \fi
    \fi
    \hep@abstract@next
  }%
  \renewcommand*\endabstract{%
    \ifhep@abstract@collecting\else\hep@abstract@end\fi
  }%
  \renewcommand*\hep@abstract@place{%
    \global\hep@title@madetrue
    \ifhep@abstract@stored
      \global\hep@abstract@storedfalse
      \begin{abstract}\hep@abstract@text\end{abstract}%
    \fi
  }
}
%    \end{macrocode}
% \end{macro}
% \end{macro}
% \end{macro}
% \end{environment}
%
% \subsection{Thanks}
%
% \begin{macro}{\thanks} Redefine thank to have a optional argument for a reference label.
% In the preamble "\thanks" without label directly after "\author", "\editor", or "\endorser" is attached to this person.
%    \begin{macrocode}
\let\hep@thanks\thanks
\AtEndDocument{\let\thanks\hep@thanks}
\DeclareRobustCommand\thanks[2][]{%
  \ifblank{#1}{%
    \hep@person@attach{\thanks{#2}}{\hep@thanks@note{}{#2}}%
  }{%
    \hep@thanks@note{#1}{#2}%
  }%
}
\newcommand\hep@thanks@note[2]{%
  \AfterEndPreamble{%
    \if\relax#1\relax%
      \footnotemark%
    \else%
      \renewcommand*\thefootnote{\textsuperscript{\@fnsymbol\c@footnote}}%
      \protect\refstepcounter{footnote}\protect\label{#1}%
      \renewcommand*\thefootnote{\@arabic\c@footnote}%
    \fi%
    \protected@xdef\@thanks{%
      \@thanks\protect\footnotetext[\the\c@footnote]{#2}%
    }%
    \if@twocolumn
      \protected@xdef\@bs@thanks{%
        \@bs@thanks\protect\footnotetext[\the\c@footnote]{#2}%
      }%
    \fi%
  }%
}
%    \end{macrocode}
% \end{macro}
%</package>
%
% \section{Test}
%
%<*testone|testtwo>
%
%     \begin{macrocode}
%<testone>\documentclass{article}
%<testtwo>\documentclass[twocolumn]{article}

%% \usepackage[lang=english]{hep-paper}
\usepackage[math]{blindtext}
\usepackage{hep-title}
\usepackage{hyperref}
\def\BackrefFootnoteTag{}\usepackage{footnotebackref}

\doi{10.1234/5678}

\arxiv{1234.5678}

\preprint{Preprint}

\series{Series}

\title{Title}

\subtitle{Subtitle}

\editor[one]{First Editor}
\editor[one,two,second]{Second Editor}
\email[second]{second@mail.com}
\editor[one,two,three]{Third Editor \email{third@mail.com}}

\author[three,four]{First Author}
\author[fourth, two,four,thanks]{Second Author}
\email[fourth]{fifth@mail.com}
\thanks[thanks]{Thanks}
\author[two,three,four]{Third Author \email{sixth@mail.com}}
\author{Fourth Author}

\endorser[four]{First Endorser}
\endorser[five]{Second Endorser}
\endorser{Third Endorser}

\affiliation[one]{First Affiliation}
\affiliation[two]{Second Affiliation}
\affiliation[three]{Third Affiliation}
\affiliation[four]{Fourth Affiliation}
\affiliation[five]{Fifth Affiliation}

\date{Date}

\keywords{Keyword 1, keyword 2, keyword 3}

\begin{document}
%<testone>
%<testone> \maketitle

%<testone>\begin{abstract}
%<testtwo>\begin{abstract*}
\blindtext
%<testone>\end{abstract}
%<testtwo>\end{abstract*}

\noindent First page footnote.\footnote{Footnote}

\blinddocument

\end{document}
%    \end{macrocode}
%
%</testone|testtwo>
%
% \section{Readme}
%
%<*readme>
%
%    \begin{macrocode}
# The `hep-title` package

Extended title data: subtitles, series, preprints, affiliations,
ORCID, keywords and abstracts

## Introduction

The `hep-title` package extends the title macros of the standard classes
with series and subtitle metadata, preprint and reprint blocks, DOI and
arXiv helpers, author, editor, endorser, affiliation, email and ORCID
helpers, date and abstract formatting, keywords, and two-column
title-aware abstracts.

To use the `hep-title` package include it with `\usepackage{hep-title}`.

## Author

Jan Hajer

## License

This file may be distributed and/or modified under the conditions of the
`LaTeX` Project Public License, either version 1.3c of this license or
(at your option) any later version. The latest version of this license is
in `http://www.latex-project.org/lppl.txt` and version 1.3c or later is
part of all distributions of LaTeX version 2005/12/01 or later.
%    \end{macrocode}
%
%</readme>
%
% \Finale

\endinput

% \PrintIndex
% makeindex -s gglo.ist -o hep-title-implementation.gls hep-title-implementation.glo
% makeindex -s gglo.ist -o hep-title-implementation.ind hep-title-implementation.idx
